\section{Fase 3: Desarrollo del motor de estados}
\label{sec:Fase3Resumen}

La tercera fase convierte el simulador en un motor de juego completo.
La Fase 2 dejó resuelta la \emph{presentación}: el reloj late, el motor
de estados escribe el estado en la lista de segmentos y el motor gráfico
compone la pantalla con repintado completo. Pero lo que se controla es
todavía una \emph{demostración}: una secuencia de prueba escrita en el
código, sin entrada, sin reglas y sin jugador. La Fase 3 construye la
parte más compleja del simulador, \textbf{el motor de estados}: la pieza
que piensa. En cada \textit{tick} lee las pulsaciones pendientes, aplica
las reglas del juego y decide qué segmentos están encendidos y dónde, si
hay colisión con el jugador aumentando o disminuyendo contadores, y
controla la vida restante hasta el fin de la partida.

Sobre el contrato fundacional de la Fase 2 —«la lista ES el estado»— el
trabajo de esta fase se apoya en una sola lista de segmentos en memoria,
un solo escritor (el motor de estados) y tantos lectores como vías de
presentación existan (hoy el \textit{canvas} del navegador; mañana el
LCD del ESP32). El motor debe ser además independiente: del director
recibe únicamente el \textit{tick}, y el resto de datos los maneja a
través del puente de datos (\emph{comunListas}).

La dificultad propia de la fase no es mover figuras: es el
\textbf{control de colisiones}, delicado por una razón de fondo del
modelo físico elegido: los segmentos de un cartucho Game \& Watch son
\emph{piezas separadas}, cada una con su sprite independiente, sin
tablero de celdas ni escena continua. En las máquinas reales la
colisión ni siquiera se calcula: el cristal líquido enciende lo que el
juego le dice, y es el \emph{firmware} original quien, en su ROM de unos
pocos kilobytes, decide con tablas y comparaciones cuándo el
paracaidista toca al tiburón. Aquí el motor tendrá que hacer ese trabajo
de forma explícita, con una decisión de diseño que recorre toda la fase:
\emph{el cartucho define, el motor ejecuta}.

El trabajo se organiza en catorce iteraciones, que se resumen a
continuación; el detalle completo de cada una está en el apéndice C.

\subsection{Iteraciones 1 y 2: El teclado}
\label{sec:F3Teclado}

\subsubsection{Primera iteración: Análisis del manejo del teclado}

\textbf{Análisis conceptual y objetivos.} El primer paso del motor es la
entrada: que una pulsación —por el botón asociado o por el teclado—
llegue al motor y mueva el segmento del jugador. Antes del código, la
iteración cierra cuatro decisiones:

\begin{itemize}
	\item \textbf{El canal: la pulsación como dato, no como evento de DOM.}
	El juego no debe conocer el DOM, igual que el firmware del ESP32 no
	conoce «botones», solo pines que se ponen a nivel. El \emph{baseDirector}
	reenvía a \texttt{juegoDirector.pulsar(numero)} el \emph{número físico}
	del botón (nada de identificadores del cartucho ni del evento), el
	director solo acepta pulsaciones con el juego en marcha y las reenvía
	tal cual al motor. La tabla número $\rightarrow$ dirección viaja en el
	código del motor, con la misma silueta que tendrá el día que las reglas
	viajen en el cartucho.
	\item \textbf{El ritmo: un buffer de movimientos pendientes.} Las
	máquinas reales muestrean la entrada al ritmo del juego. Las
	pulsaciones se encolan en un buffer FIFO con tapa (8 entradas) y el
	motor las procesa al empezar el ciclo siguiente, en un paso (0) previo
	a todo lo demás: leer pines al inicio del frame, computar, presentar.
	El estado nunca cambia a mitad de presentación, y el orden del
	\textit{tick} queda sagrado: (0) entrada, (1) apagar banderas de
	cambio, (2) secuencia.
	\item \textbf{La combinación de teclas es una propiedad del cartucho.}
	En el Donkey Kong de bolsillo, ARRIBA+DERECHA no son dos movimientos:
	es un salto diagonal. Que un juego permita combinar dos teclas lo
	decide el diseñador, declarado como un indicador más del JSON
	(\texttt{combinar}): con combinación activa el motor vacía todo el
	buffer del tick y lo reduce a los ejes, moviendo el jugador una sola
	vez con el vector resultante; sin ella, cada entrada es un paso
	independiente; direcciones opuestas en la misma ventana se anulan.
	\item \textbf{Qué máquina ejecuta qué cartucho: setup $\geq$ cartucho.}
	El setup declara la máquina física y el cartucho usa un subconjunto de
	sus botones: se puede jugar a un juego de dos botones en una máquina de
	cuatro, pero no al contrario. Si el cartucho declara un botón que la
	máquina no tiene, el cartucho NO carga (error bloqueante en la barra
	de estado); los botones de la máquina que el juego no usa se eliminan
	de la botonera.
\end{itemize}

\subsubsection{Segunda iteración: El motor responde al teclado}

\textbf{Decisiones de diseño y arquitectura.} La iteración anterior
deja las decisiones cerradas; esta las hace carne. El motor aprende a
responder: una pulsación encola, el \textit{tick} la consume y el
jugador se mueve. Moverse es escribir en la lista: un índice de
segmento jugador y una tabla \texttt{MOVER} número $\rightarrow$ vector,
escritos en el código del motor como lo estaba la secuencia de prueba;
al \texttt{iniciar()} el motor guarda las posiciones originales de la
lista y mover es escribir \emph{x}/\emph{y} y marcar la bandera de
cambio; relanzar restaura las posiciones de origen en el mismo barrido
que limpia banderas. El jugador no se mueve libremente: solo puede
seguir los enlaces de izquierda, derecha, arriba y abajo, por lo que no
puede caerse por los bordes. La combinación viaja por el puente:
\emph{comunListas} gana una casilla de configuración del juego que la
autoría escribe al cargar el cartucho y el motor lee en cada
\textit{tick}. El motor gráfico no se toca: la composición pura
presenta lo que la lista dice, y la secuencia de fondo sigue latiendo
en paralelo con el movimiento del jugador, compartiendo el mismo
\textit{tick} y la misma lista.

\textbf{Verificación.} La simulación ejercita el contrato completo:
pulsación con el juego parado no se encola; con \texttt{combinar}
activo, dos pulsaciones en la misma ventana dan UN paso diagonal
verificado a píxel; sin combinar, dos pasos en orden; opuestas no
mueven; la tapa descarta el exceso con aviso; relanzar restaura la
posición de origen; un cartucho que pida un botón inexistente NO
carga. Cero errores de consola con GW\_DEBUG apagada.

\subsection{Tercera iteración: La marca de primer jugador}
\label{sec:F3PrimerJugador}

\textbf{Análisis conceptual y objetivos.} El motor tomaba como jugador
activo el primer segmento de tipo jugador que encontraba en la lista:
una búsqueda que no decía por dónde empieza la partida. Esta iteración
retrocede un paso al diseñador para dársela: una fila «Jugador
inicial» en el panel de propiedades, visible solo para los segmentos
de tipo jugador.

\textbf{Decisiones de diseño y arquitectura.} La marca vive como la
propiedad \texttt{primero} del componente: entra en el historial como
cualquier otro campo y viaja en el GWI como
\texttt{primero: true}. Marcar dos segmentos es una acción atómica del
historial: se quita la marca al anterior y se da al nuevo en el mismo
bloque, de modo que deshacer revierte ambas cosas de un golpe; solo
puede haber una marca y prevalece la última intención del autor.
Cambiar de tipo descarta la marca, igual que los enlaces de dirección.
El motor solo cambia una prioridad: \texttt{buscarJugador} consulta
primero la marca del diseñador y, si no hay, cae a la búsqueda
anterior, que queda como salvavidas: el sistema funciona siempre,
marcado o no.

\subsection{Cuarta iteración: El generador de aleatorios}
\label{sec:F3PRNG}

\textbf{Análisis conceptual y objetivos.} Los G\&W reales generan
números aleatorios con un contador que apunta a una zona de la ROM, sin
tablas dedicadas. En el simulador no hay ROM que leer, así que se
sustituye por un generador matemático que produce una secuencia
determinista a partir de una semilla: misma semilla, misma secuencia,
misma partida, condición indispensable para depurar y para el
hardware.

\textbf{Decisiones de diseño y arquitectura.} Se adopta un generador
lineal congruencial (LCG) de 16 bits, la opción clásica para hardware
pequeño: avance $\text{estado} = (a \cdot \text{estado} + c)
\bmod 2^{16}$ con las constantes clásicas $a = 25173$ y $c = 13849$, de
la familia usada por los BASIC de 8 bits; el desbordamiento natural del
\texttt{uint16\_t} es el módulo, con coste cero en el ESP32. La lectura
se hace con $\text{estado} \bmod N$, suficiente para los $N$ del juego
(rutas, dígitos). La ubicación es el \emph{director del juego}: aunque el
consumidor inmediato es el motor, el director puede necesitar aleatorios
para eventos propios, y el principio del proyecto es que el director es
el coordinador al que el motor consulta. La API es una única puerta,
\texttt{aleatorio(max)}; la semilla es una constante editable del
fichero y \texttt{arrancar()} reinicia el estado del LCG a la semilla,
de modo que cada partida empieza con la misma secuencia, igual que en
las máquinas reales. El diagnóstico va bajo GW\_DEBUG.

\textbf{Verificación.} El director inicializa el sistema, pide diez
números seguidos, reinicializa y vuelve a pedirlos: ambas tandas son
idénticas, verificando el determinismo de punta a punta.

\subsection{Quinta iteración: Lanzar las rutas}
\label{sec:F3Rutas}

\textbf{Análisis conceptual y objetivos.} Con el generador verificado,
esta iteración le da su primer uso real: elegir rutas. La secuencia de
prueba desaparece y, por primera vez, algo definido en el cartucho
decide lo que se ve en pantalla. Una ruta es la cadena que ya viaja en
el GWI: un segmento \textbf{inicio} (numerado con el campo
\texttt{ruta}) y los segmentos que cuelgan de sus enlaces
\texttt{siguiente}. El «elemento» que recorre la ruta no es una entidad
nueva: es el segmento encendido ahora en esa cadena — recorrer una ruta
es apagar el segmento actual y encender el que marca su
\texttt{siguiente}, el mismo mecanismo del jugador, solo que con el
reloj.

\textbf{Decisiones de diseño y arquitectura.} El paso de avance del
\textit{tick} cambia: en el barrido solo participan los segmentos de
tipo juego e inicio (función \texttt{esDeRuta}); antes de tocar nada se
toma una FOTO de los segmentos de ruta visibles, de modo que cada uno
avanza exactamente un paso aunque su siguiente también estuviera
visible; si el actual no tiene siguiente, la ruta muere y el segmento
final se apaga sin rastro. La política de lanzamiento, de momento de
pruebas, encadena rutas sin huecos: mientras una ruta vive no nace otra,
y en el mismo \textit{tick} en que muere el PRNG elige la siguiente, que
nace al instante — la partida nunca muestra un \textit{tick} vacío. El
motor no añade estado propio: si una ruta vive, es porque su segmento
actual tiene la bandera visible encendida; la lista sigue siendo el
estado. El mecanismo no limita las rutas simultáneas (varias a la vez
estilo Game B quedará como parámetro del cartucho en fases posteriores):
es la política de lanzamiento la que impone el encadenado, no el
barrido.

\textbf{Verificación.} Las tres rutas del cartucho de prueba se lanzan
de forma aleatoria hasta el final, recorriendo la cadena completa
verificada segmento a segmento; con la misma semilla, dos partidas
eligen el MISMO orden de rutas; el jugador sigue moviéndose por sus
enlaces mientras las rutas corren.

\subsection{Iteraciones 6 a 8: Las colisiones}
\label{sec:F3Colisiones}

\subsubsection{Sexta iteración: Análisis del problema de las colisiones}

\textbf{Análisis conceptual y objetivos.} Jugador y rutas conviven en
pantalla sin interferir. La decisión de diseño es que la colisión
tampoco sea geometría: \textbf{la colisión es un enlace declarado}. En
un segmento de juego, la autoría marca con qué segmento del jugador
colisiona; en tiempo de ejecución, el motor solo mira una cosa: ¿está
el jugador visible en el lugar que marca el enlace? Sin cálculo de
cajas, sin solapes, sin umbrales: una comparación de índices, como en
las máquinas reales, donde el programador decidía qué solapes
contaban. Un punto de colisión es un punto de decisión con dos caminos:
\emph{colisión positiva} (el jugador está visible) y \emph{colisión
	negativa} (no lo está). Cada desenlace puede hacer tres cosas definidas
por la autoría: seguir una cadena (los enlaces ya existentes
\texttt{salvado} y \texttt{muerte} reciben su semántica definitiva:
desenlace con jugador y sin jugador), actualizar un contador interno
una sola vez, en el instante de la decisión — con un indicador
\emph{vida}, \emph{muerte} o \emph{nada} que viaja junto a su enlace —,
o nada: lo que no está declarado, no pasa; ninguna colisión puede
quedar sin comportamiento. Distinción importantísima cerrada aquí: el
\emph{contador interno} que el motor mantiene no es lo mismo que los
segmentos de tipo vida o muerte, que son su \emph{presentación} en
pantalla.

\subsubsection{Séptima iteración: Las colisiones en el diseñador}

\textbf{Decisiones de diseño y arquitectura.} Los segmentos de juego
ganan la marca de colisión — elegir el segmento jugador con el que
chocan — y los contadores de desenlace (\emph{vida}, \emph{muerte} o
\emph{ninguno}). Las propiedades nuevas entran en el historial como
cualquier otro campo, sobreviven a las uniones re-apuntándose y se
reproducen al recargar el GWI. La verificación calcula el inverso de
la colisión (el jugador sabe quién choca contra él, \texttt{colisionDe}),
avisa de un jugador sin puntos de colisión y bloquea una colisión cuyo
destino no es un jugador, una cadena de desenlace que enlaza con el
jugador, o un jugador inalcanzable por los enlaces de dirección: no se
permiten datos incoherentes. La regla «un jugador, un punto de
colisión» se sostiene (la multi-colisión queda anotada como ampliación),
y el mapa pinta los puntos de colisión en verde azulado con su trazo
al jugador, para verlos de un vistazo.

\subsubsection{Octava iteración: El motor controla las colisiones}

\textbf{Decisiones de diseño y arquitectura.} El paso (2) del
\textit{tick} gana una regla que va ANTES del avance: un segmento
visible de ruta que declara colisión NO avanza por \texttt{siguiente} —
decide. La comprobación es una lectura de bandera del segmento jugador
señalado; si está visible hay colisión positiva y la ruta sigue
\texttt{salvado}; si no, negativa y sigue \texttt{muerte}; sin enlace,
el elemento desaparece (regla por defecto). La cadena elegida se
comporta como cualquier ruta: sin estado de colisión aparte, las
banderas son el estado, y la foto del barrido evita dobles decisiones.
El contador indicado se toca UNA sola vez, en la decisión, nunca en
cada \textit{tick} de la animación del desenlace. El momento de la
decisión es el \textit{tick} en que el punto de colisión ESTÁ visible,
sin un tick de espera: si el jugador llega tarde, es negativa; el
\emph{timing} es parte del juego, como en las máquinas reales. Los
contadores internos se muestran en la barra de diagnóstico
(Tick N · Vidas: V · Muertes: M) y \texttt{iniciar()} los restaura
(vidas = número de segmentos de tipo vida del cartucho, para que
coincida con su presentación; muertes = 0).

\textbf{Problemas detectados y corrección.} En las pruebas se observa
que con solo tres rutas de salida, la misma ruta se repite demasiadas
veces seguidas y la partida parece un bucle. La corrección no toca el
generador: el motor guarda una lista con los números de las últimas $n$
rutas lanzadas ($n$ constante ajustable, inicialmente 2) y descarta la
ruta elegida solo si las últimas $n$ son todas ella; con $n = 2$ una
ruta puede salir dos veces seguidas, nunca tres. La elección sigue
siendo función del estado del generador: libre y determinista (misma
semilla = misma partida), y al iniciar la lista se vacía.

\subsection{Iteraciones 9 y 10: Los contadores en pantalla}
\label{sec:F3Contadores}

\subsubsection{Novena iteración: un callejón sin salida}

\textbf{Problema detectado.} Los contadores internos vivían solo en la
barra de depuración; la iteración se propone presentarlos en pantalla
como en las máquinas reales, donde el arte lleva un «8» impreso en cada
posición del número y el sistema sabe qué barras enciende cada dígito
(un 1 es \emph{b+c}, un 2 es \emph{a,b,g,e,d}; el «8» puede dibujarse
normal o en cursiva, con tal de que las siete barras se extraigan sin
ambigüedad). La implementación se apoyó en la \emph{unión} de segmentos
existente y lo estiró para cubrir la jerarquía contador
$\rightarrow$ dígitos $\rightarrow$ barras. Los intentos fueron varios
y ninguno cerró: reciclar una barra como dígito dejaba semántica
torcida que la verificación rechazaba; separar «segmentos ficticios»
conviviendo con la unión clásica fue frágil — piezas agrupadas que no
respondían al clic, «Separar» que no separaba, errores de guardado — y
cada arreglo tapaba un agujero abriendo otro.

\textbf{La decisión.} El problema no eran los contadores: era el
sistema de unión. Fusionar los segmentos en uno de los existentes (el
absorbido desaparece) no podía sostener encima una jerarquía de dos
niveles con piezas que nacen y se separan. La ventaja de las
iteraciones pequeñas es que se puede parar y retroceder antes de
continuar: la iteración se declara fallida y se rehace el mecanismo de
unión de arriba abajo.

\subsubsection{Décima iteración: un nuevo sistema de unión y los dígitos}

\textbf{Decisiones de diseño y arquitectura.} La unión deja de fundir
los segmentos en uno de los existentes: al unir \texttt{seg\_020} y
\texttt{seg\_030} se crea una pieza nueva (\texttt{seg\_F01}) con una
imagen nueva (\texttt{img\_F01}, fusión de ambas) que absorbe a los
originales, hereda sus propiedades y re-apunta sus enlaces externos;
los originales siguen existiendo en el cartucho, marcados
\texttt{visible:false}, sin presentarse ni manejarse. «Unir», «Hacer
dígito» y «Hacer contador» pasan a ser el mismo mecanismo de
absorción, y «Separar» su operación inversa, siempre deshacible.

Sobre esa base se montan los contadores: el autor selecciona las siete
barras de un «8» y el sistema clasifica cada barra por su geometría
(las tres horizontales por altura son \emph{a}, \emph{g}, \emph{d}; las
verticales de la derecha \emph{b} y \emph{c}, las de la izquierda
\emph{f} y \emph{e}), creando un segmento de tipo \textbf{dígito} que
agrupa las barras mediante los enlaces \texttt{barra\_a} a
\texttt{barra\_g}; con los dígitos marcados (hasta 4) se crea el
\textbf{contador} con los enlaces \texttt{digito\_1} (las unidades,
obligatorio) a \texttt{digito\_4}. Las barras y los dígitos dejan de ser
accesibles en el editor (se selecciona la pieza completa), pero siguen
viajando individualmente en el cartucho: el motor enciende las barras
reales de cada posición. El contador declara su número de contador
interno (1 a 4, autoasignado), su valor inicial, si la cuenta es
ascendente o descendente (que no baja de cero) y si presenta o no el
cero. Los valores menores que $10^n$ van rellenos con blancos a la
izquierda si lo necesitan, y los mayores o iguales que $10^n$ se
presentan rellenos con ceros: el contador interno sigue creciendo y la
presentación muestra los últimos $n$ dígitos.

Como consecuencia, la funcionalidad «Posicionar juntos» de la autoría
—nacida precisamente para simular los contadores numéricos de las
vidas— se retira del editor: con los contadores reales ya no tiene
sentido.

\subsection{Undécima iteración: El fin de la partida}
\label{sec:F3Fin}

\textbf{Análisis conceptual y objetivos.} Una partida no terminaba
nunca. Falta definir las características generales del juego que no
pertenecen a ningún segmento de arte, y que el motor reaccione a ellas.
Como el GWI no se monta a mano sino con el editor, se crea un tipo nuevo
de segmento, \textbf{cartucho} (tipo 7): una pieza ficticia que nunca
se pinta y cuya ficha se abre con un clic en cualquier punto vacío del
mapa. Hay exactamente una por cartucho, creada siempre automáticamente,
que el autor no puede borrar ni duplicar; en el GWI no viaja como
segmento, sino como propiedad del cartucho al nivel superior:

\begin{lstlisting}[language=json]
	{ "fin": { "mecanismo": "vida" } }
	{ "fin": { "mecanismo": "contador", "contador": 1 } }
\end{lstlisting}

\textbf{Decisiones de diseño y arquitectura.} El autor elige uno de dos
mecanismos, que no pueden convivir:

\begin{itemize}
	\item \textbf{Por contador:} uno de los cuatro contadores internos,
	descendente, iniciado al número de vidas; la partida termina cuando el
	contador \textbf{transiciona de mayor que cero a cero}. Empezar en cero
	(que significa «nunca termina») es error bloqueante al guardar, igual
	que no tener contador en pantalla.
	\item \textbf{Por iconos de vida o muerte:} los segmentos de muerte
	(los tres tiburones bajo el aviso) nacen apagados y cada muerte enciende
	el último no visible; la partida termina al encenderse el último. Los
	de vida nacen todos encendidos y cada pérdida apaga el último; termina
	cuando no queda ninguno visible. Las vidas iniciales no hay que
	definirlas: salen contando los iconos de la familia. Solo puede haber
	una familia de vida o una de muerte por cartucho.
\end{itemize}

La asociación del fin es automática cuando ayuda: al convertir un
segmento a tipo vida o muerte, si el fin está sin definir, se asocia a
esa familia con su aviso; si el segmento era el último de su familia,
el fin se retira. Asociación y retirada viajan en la misma acción
atómica del historial que el cambio de tipo. Para que las dos etapas
hablen el mismo idioma, la enumeración de tipos de segmento y de
mecanismos de fin vive en el puente (\emph{comunListas}): comparar
contra constantes con nombre en vez de dígitos mágicos.

En el motor, el estado \textbf{terminada} congela la partida como en
las máquinas reales: la pantalla queda tal cual estaba, el barrido de
rutas y la entrada del jugador dejan de avanzar y el reloj se detiene,
con el aviso «Partida terminada — pulsa Start»; el Start existente
reinicia semilla, cola, contadores e \texttt{iniciar()}: la nueva
partida es el arranque de siempre, sin piezas nuevas.

\textbf{Problemas detectados y corrección.} En la verificación por
iconos de muerte, la primera muerte fallaba porque el tiburón que
quedaba fijo era el primero en lugar del último: el orden de encendido
de la familia debe marcarlo la cadena de ruta, no la numeración del
editor. Tras ajustarlo, los tres mecanismos (contador, iconos de
muerte, iconos de vida) terminan la partida correctamente.

\subsection{Duodécima iteración: Cambios menores, elementos sin interacción y la Secuencia}
\label{sec:F3Menores}

\textbf{Análisis conceptual y objetivos.} El sistema ya funciona y se
puede jugar; quedan los temas menores de las máquinas reales: nuevas
validaciones al guardar, elementos que se presentan pero no
intervienen, y la congelación del juego durante la secuencia de muerte.

\begin{itemize}
	\item \textbf{Validaciones adicionales:} el fin por vida o muerte exige
	que exista al menos un segmento de ese tipo; todo cartucho necesita al
	menos un jugador, y todo jugador algún enlace de dirección (un jugador
	que no puede moverse nunca es correcto); y los enlaces de dirección
	del jugador calculan su inverso automáticamente: al enlazar
	\texttt{A $\rightarrow$ izquierda: B}, en B se escribe
	\texttt{derecha: A}, y lo mismo al quitar, en la misma acción del
	historial.
	\item \textbf{Elementos sin interacción:} dos tipos nuevos. El
	\textbf{adicional} se presenta pero no interviene: sin enlace
	\texttt{siguiente} nace encendido y queda fijo para siempre (el
	helicóptero); con enlace nace apagado y solo lo enciende su propio
	bucle. El \textbf{adicional de inicio} lanza la secuencia y nace
	encendido, alternando un segmento por tick, indefinidamente (las
	aspas). Sus enlaces \texttt{siguiente} solo pueden apuntar a
	segmentos de estos tipos.
	\item \textbf{La Secuencia:} animaciones de fondo que no forman parte
	del juego (las aletas de tiburón bajo las barcas). Cada cartucho puede
	disponer de una, definida por sus \textbf{ticks} de espera y una lista
	de hasta cinco segmentos lanzados en orden, uno por tick; los
	segmentos pueden ser de juego, inicio o adicional, y pueden pertenecer
	a la vez a otras secuencias (las aletas merodean solas y, a la vez, son
	eslabones de la secuencia de muerte). El ciclo real dura $n + m$
	\textit{ticks} ($n$ de espera más $m$ los que dura la animación).
	\item \textbf{La congelación por desenlace:} cuando una colisión
	lanza una cadena de vida o de muerte, en la máquina real el resto del
	juego se congela hasta que termina la animación. Una bandera nace con
	la cadena, se propaga por sus enlaces \texttt{siguiente} y cae al morir
	el último eslabón; mientras está encendida, las rutas no avanzan ni
	nacen y el jugador no responde — solo sigue vivo el bucle de los
	elementos adicionales (las aspas). Una muerte instantánea sin animación
	no congela nada; y como salvaguarda, si una cadena marcada se rompe por
	un enlace colgante, la marca se limpia sola: un error de autoría no
	puede congelar el juego para siempre.
\end{itemize}

\textbf{Problemas detectados y corrección.} Esta iteración obligó a
depurar tres errores con síntomas parecidos y causas distintas, un buen
ejemplo del valor del diagnóstico:

\begin{enumerate}
	\item \textbf{Interrupción en cada colisión:} el diseño original ponía
	el contador de la Secuencia a cero ante cualquier colisión decidida. En
	el cartucho real, donde las colisiones son constantes, la animación
	nunca completaba un ciclo. La corrección limita la suspensión al ciclo
	de muerte: las colisiones positivas y la congelación por desenlace
	salvado no la afectan.
	\item \textbf{Puente de datos desactualizado:} el editor solo escribía
	la Secuencia en el puente al cargar el cartucho, de modo que definirla
	en la ficha y arrancar sin recargar dejaba al motor sin verla. La
	corrección sincroniza el puente al final de cada reconstrucción de la
	ficha: cualquier edición queda visible al instante.
	\item \textbf{Ciclo de muerte atascado:} los segmentos de la Secuencia
	pueden ser a la vez eslabones de las cadenas de desenlace (las aletas
	compartidas), y el barrido los excluía del avance por estar encendidos
	por la animación, con lo que la cadena de muerte se detenía en ellos
	para siempre y no se producían muertes. La corrección distingue los
	dos contextos: solo el barrido de rutas respeta la marca de animado; un
	eslabón de la cadena de desenlace avanza aunque esté encendido por la
	Secuencia. Como remate, al nacer el ciclo de muerte el segmento activo
	de la animación se apaga y la Secuencia se reinicia por completo.
\end{enumerate}

La localización se apoyó en la barra de diagnóstico, que con GW\_DEBUG
muestra en todo momento el estado de la Secuencia (contador de
\textit{ticks} pendientes, segmento en curso y si anima o está
suspendida), permitiendo distinguir si el fallo estaba en el motor, en
el puente de datos o en la presentación.

\subsection{Decimotercera iteración: Segmentos que se presentan juntos}
\label{sec:F3Juntos}

\textbf{Análisis conceptual y objetivos.} En la máquina real hay
parejas de segmentos que se presentan siempre a la vez al recorrer una
secuencia (un nadador perseguido por una aleta), pero cuyo componente
secundario también aparece solo en la animación de fondo. Ni la unión
(que fusiona para siempre) ni los enlaces lo cubren.

\textbf{Decisiones de diseño y arquitectura.} El segmento principal
gana el parámetro «Mostrar junto a», que enlaza con el secundario; en
el secundario, «Mostrado junto a» informa en automático de su
principal. Cuando se enciende o apaga el principal se hace siempre lo
mismo con el secundario; pero un secundario encendido directamente (por
la animación) no altera al principal. El secundario no puede formar
parte de ninguna red de enlaces, no reporta colisión y solo puede ser
de tipo juego, todo verificación bloqueante al guardar.

\textbf{Problema detectado y corrección.} En la verificación se detectó
que al lanzarse una ruta de muerte durante una animación, esta
continuaba por donde estaba al finalizar la cadena, cuando debía
cortarse por completo y recomenzar el conteo de \textit{ticks}. Corregido
y reverificado, las aletas funcionan en solitario durante las
animaciones y persiguiendo a los paracaidistas en las secuencias de
muerte.

\subsection{Cierre de la fase: Batería con las pruebas efectuadas y su resultado}
\label{sec:F3Pruebas}

Como en las fases anteriores, el cierre se certifica con una iteración
de validación que recorre todos los cambios incorporados, desde la
carga del cartucho hasta la presentación en pantalla con el GWI real
del Parachute:

\begin{longtable}{|c|l|p{8.2cm}|c|}
	\hline
	\textbf{Nº} & \textbf{Grupo} & \textbf{Prueba} & \textbf{Resultado} \\
	\hline
	\endfirsthead
	\hline
	\textbf{Nº} & \textbf{Grupo} & \textbf{Prueba} & \textbf{Resultado} \\
	\hline
	\endhead
	
	1 & Jugador & El jugador queda visible y presentado en su posición
	de inicio al primer \texttt{iniciar}, sin necesidad de pulsar nada. &
	OK \\
	\hline
	2 & Jugador & Los botones mueven al jugador por sus enlaces de
	dirección; los enlaces inversos se mantienen automáticamente al marcarlos
	y al quitarlos. & OK \\
	\hline
	3 & Jugador & La verificación al guardar bloquea un cartucho con un
	jugador sin ningún enlace de dirección a otro jugador. & OK \\
	\hline
	4 & Rutas & Los segmentos de juego nacen por azar, avanzan por sus
	enlaces un segmento por tick y mueren al llegar al final de su
	concatenación. & OK \\
	\hline
	5 & Azar & Con la misma semilla, dos partidas son idénticas: mismas
	rutas naciendo en los mismos ticks y mismos desenlaces. & OK \\
	\hline
	6 & Colisiones & Una colisión solo se decide entre el jugador y un
	segmento visible que reporta colisión; el desenlace salvado enciende su
	cadena y el mortal la suya. & OK \\
	\hline
	7 & Congelación & Mientras corre una cadena de desenlace el juego queda
	congelado: no nacen rutas y el jugador no responde; al morir el último
	eslabón el juego se reanuda. & OK \\
	\hline
	8 & Contadores & La familia de vida se enciende desde el primero y la
	de muerte acumulativamente desde el último, en el orden definido por su
	cadena y no por la numeración del diseñador. & OK \\
	\hline
	9 & Contadores & Los contadores de dígitos presentan su valor con
	segmentos de siete segmentos y lo actualizan en el mismo tick del
	cambio. & OK \\
	\hline
	10 & Fin & La partida termina al agotarse un contador, al encenderse
	todos los iconos de muerte o al perderse todas las vidas, según el
	mecanismo declarado en el cartucho. & OK \\
	\hline
	11 & Adicionales & Un segmento \emph{adicional} sin siguiente queda
	fijo encendido (helicóptero); con siguiente avanza en su propio bucle,
	y un \emph{inicioAdicional} nace encendido y alterna solo (aspas), sin
	que nada los interrumpa. & OK \\
	\hline
	12 & Secuencia & La Secuencia lanza sus segmentos en orden, uno por
	tick, y se relanza desde el primero tras esperar los ticks declarados
	(ciclo completo de $n + m$ ticks). & OK \\
	\hline
	13 & Secuencia & Las colisiones y la congelación por desenlace salvado
	NO interrumpen la Secuencia; solo el nacimiento de un ciclo de muerte
	la interrumpe. & OK \\
	\hline
	14 & Secuencia & Al nacer el ciclo de muerte, el segmento activo de la
	Secuencia desaparece, la animación se reinicia por completo y espera los
	ticks necesarios para relanzarse desde su primer segmento. & OK \\
	\hline
	15 & Secuencia & La cadena de muerte atraviesa un segmento encendido
	por la Secuencia (las aletas compartidas): lo toma como eslabón, lo apaga
	y sigue, produciendo la muerte. & OK \\
	\hline
	16 & Acompañantes & Un secundario «Mostrado junto a» se enciende y
	apaga siempre con su principal, incluidas las cadenas de desenlace, y
	puede presentarse solo en la animación de fondo. & OK \\
	\hline
	17 & Acompañantes & La verificación al guardar bloquea un secundario
	con cualquier enlace: con siguiente, apuntado por otro, con colisión,
	de tipo distinto a juego o con dos principales. & OK \\
	\hline
	18 & Depuración & La barra de estado con GW\_DEBUG muestra el estado de
	la Secuencia (contador pendiente, segmento en curso y si anima o está
	suspendida). & OK \\
	\hline
	19 & Integración & Partida completa con el GWI real del Parachute:
	tres ciclos de muerte que atraviesan las aletas encendidas por la
	animación, tres iconos de muerte y partida terminada por fin de
	muerte. & OK \\
	\hline
	\caption{Batería de pruebas de validación de la Fase 3.}
\end{longtable}

Con esta batería en verde, la Fase 3 se cierra con el simulador
jugable de punta a punta: el cartucho define y el motor ejecuta;
la misma disciplina con la que afrontaremos la Fase 4, donde los
parámetros que gobiernan el ritmo de la partida saldrán del código
para viajar también en el cartucho.

\clearpage